%% annex/p02/propositional-metatheory-supplement.tex --- retained editorial proofs
\noindent\textit{Reader note.} These three proofs support only the separately
attributed duality, interpolation, and dual-Horn additions retained after the exact
source chapter.  The exact inherited proofs remain inline in their governed
source scopes; this supplement neither duplicates nor replaces them.

\begin{annexitem}{annex:p2.1.8}{thm:p02-propositional-duality}{Duality theorem}
\begin{proof}
Throughout, \(\bar v\) is the valuation with \(\bar v(p_n)=\mathbf T\) iff
\(v(p_n)=\mathbf F\), for every \(n\); under the reading
\(\mathbf F=0\), \(\mathbf T=1\) this is the valuation \(1-v(p)\) of the
statement.  Complementing twice returns the original valuation:
\(\bar{\bar v}=v\).  By
Lemma~\ref{thm:p02-w02-p02-010-nnf-presentation}, every formula in negation
normal form has exactly one NNF presentation, built by the clauses
\(L\), \(\bottom\), \(\topv\), \((\beta\land\gamma)\), \((\beta\lor\gamma)\)
of Definition~\ref{def:p02-nnf-dual} with \(\beta,\gamma\) again NNF
presentations; so induction on NNF presentations covers six cases, a literal
being either a variable \(p_n\) or a negated variable \(\neg p_n\).

\emph{1. The semantic identity.}  We prove, for every NNF presentation
\(\alpha\), that every valuation \(v\) satisfies
\(v\entails\alpha^{\mathrm d}\) iff \(\bar v\not\entails\alpha\).
\begin{enumerate}
  \item \(\alpha=p_n\).  Then \(\alpha^{\mathrm d}=p_n\), and
        \(v\entails p_n\) iff \(v(p_n)=\mathbf T\) iff
        \(\bar v(p_n)=\mathbf F\) iff \(\bar v\not\entails p_n\).
  \item \(\alpha=\neg p_n\).  Then \(\alpha^{\mathrm d}=\neg p_n\), and
        \(v\entails\neg p_n\) iff \(v(p_n)=\mathbf F\) iff
        \(\bar v(p_n)=\mathbf T\) iff \(\bar v\entails p_n\) iff
        \(\bar v\not\entails\neg p_n\), using the negation clause at the first
        and last steps.
  \item \(\alpha=\bottom\).  Then \(\alpha^{\mathrm d}=\topv=\neg\bottom\).
        Since \(\widehat v(\bottom)=\mathbf F\), the negation clause gives
        \(v\entails\neg\bottom\); and \(\bar v\not\entails\bottom\).  Both
        sides hold for every \(v\).
  \item \(\alpha=\topv=\neg\bottom\).  Then \(\alpha^{\mathrm d}=\bottom\),
        so \(v\not\entails\alpha^{\mathrm d}\); and \(\bar v\entails\neg\bottom\)
        as in the previous case, so \(\bar v\not\entails\alpha\) is false.  Both
        sides fail for every \(v\).
  \item \(\alpha=(\beta\land\gamma)\).  Then
        \(\alpha^{\mathrm d}=(\beta^{\mathrm d}\lor\gamma^{\mathrm d})\), and
        \[
         v\entails\alpha^{\mathrm d}
         \ \text{iff}\
         v\entails\beta^{\mathrm d}\text{ or }v\entails\gamma^{\mathrm d}
         \ \text{iff}\
         \bar v\not\entails\beta\text{ or }\bar v\not\entails\gamma
         \ \text{iff}\
         \bar v\not\entails(\beta\land\gamma),
        \]
        by the truth clause for \(\lor\), the induction hypotheses for
        \(\beta\) and \(\gamma\), and the truth clause for \(\land\).
  \item \(\alpha=(\beta\lor\gamma)\).  Then
        \(\alpha^{\mathrm d}=(\beta^{\mathrm d}\land\gamma^{\mathrm d})\), and
        \[
         v\entails\alpha^{\mathrm d}
         \ \text{iff}\
         v\entails\beta^{\mathrm d}\text{ and }v\entails\gamma^{\mathrm d}
         \ \text{iff}\
         \bar v\not\entails\beta\text{ and }\bar v\not\entails\gamma
         \ \text{iff}\
         \bar v\not\entails(\beta\lor\gamma),
        \]
        by the truth clause for \(\land\), the induction hypotheses for
        \(\beta\) and \(\gamma\), and the truth clause for \(\lor\).
\end{enumerate}

\emph{2. Involution.}  The clauses send leaves to leaves and binary
presentations to binary presentations, so \(\alpha^{\mathrm d}\) is again an
NNF presentation and \((\alpha^{\mathrm d})^{\mathrm d}\) is defined.  We prove
\((\alpha^{\mathrm d})^{\mathrm d}=\alpha\) by induction on \(\alpha\).
If \(\alpha\) is a literal \(L\), then
\((L^{\mathrm d})^{\mathrm d}=L^{\mathrm d}=L\).  If \(\alpha=\bottom\), then
\((\bottom^{\mathrm d})^{\mathrm d}=\topv^{\mathrm d}=\bottom\); if
\(\alpha=\topv\), then
\((\topv^{\mathrm d})^{\mathrm d}=\bottom^{\mathrm d}=\topv\).
If \(\alpha=(\beta\land\gamma)\), then
\[
 ((\beta\land\gamma)^{\mathrm d})^{\mathrm d}
 =(\beta^{\mathrm d}\lor\gamma^{\mathrm d})^{\mathrm d}
 =((\beta^{\mathrm d})^{\mathrm d}\land(\gamma^{\mathrm d})^{\mathrm d})
 =(\beta\land\gamma),
\]
by the \(\land\)-clause, then the \(\lor\)-clause, then the induction
hypotheses.  If \(\alpha=(\beta\lor\gamma)\), then
\[
 ((\beta\lor\gamma)^{\mathrm d})^{\mathrm d}
 =(\beta^{\mathrm d}\land\gamma^{\mathrm d})^{\mathrm d}
 =((\beta^{\mathrm d})^{\mathrm d}\lor(\gamma^{\mathrm d})^{\mathrm d})
 =(\beta\lor\gamma),
\]
by the \(\lor\)-clause, then the \(\land\)-clause, then the induction
hypotheses.  Hence dualization maps the NNF presentations onto themselves and
is its own inverse.

\emph{3. Logical equivalence.}  Suppose \(\alpha\equiv\beta\) and let \(v\) be
any valuation.  By part~1 applied to \(\alpha\) and to \(\beta\),
\(v\entails\alpha^{\mathrm d}\) iff \(\bar v\not\entails\alpha\) iff
\(\bar v\not\entails\beta\) iff \(v\entails\beta^{\mathrm d}\), the middle step
because \(\alpha\) and \(\beta\) take the same value under \(\bar v\).  So
\(\alpha^{\mathrm d}\equiv\beta^{\mathrm d}\).  Conversely, if
\(\alpha^{\mathrm d}\equiv\beta^{\mathrm d}\), the implication just proved,
applied to the NNF presentations \(\alpha^{\mathrm d}\) and
\(\beta^{\mathrm d}\), gives
\((\alpha^{\mathrm d})^{\mathrm d}\equiv(\beta^{\mathrm d})^{\mathrm d}\), which
is \(\alpha\equiv\beta\) by part~2.

\emph{4. Valid and unsatisfiable formulas.}  All four implications use part~1
and \(\bar{\bar u}=u\).
\begin{enumerate}
  \item If \(\alpha\) is valid, then for every \(v\) the valuation \(\bar v\)
        satisfies \(\alpha\), so \(v\not\entails\alpha^{\mathrm d}\); hence
        \(\alpha^{\mathrm d}\) is unsatisfiable.
  \item If \(\alpha^{\mathrm d}\) is unsatisfiable, let \(u\) be any valuation
        and put \(v=\bar u\).  Then \(v\not\entails\alpha^{\mathrm d}\), so
        \(\bar v\entails\alpha\), that is, \(u\entails\alpha\).  Hence
        \(\alpha\) is valid.
  \item If \(\alpha\) is unsatisfiable, then for every \(v\) the valuation
        \(\bar v\) falsifies \(\alpha\), so \(v\entails\alpha^{\mathrm d}\);
        hence \(\alpha^{\mathrm d}\) is valid.
  \item If \(\alpha^{\mathrm d}\) is valid, let \(u\) be any valuation and put
        \(v=\bar u\).  Then \(v\entails\alpha^{\mathrm d}\), so
        \(\bar v\not\entails\alpha\), that is, \(u\not\entails\alpha\).  Hence
        \(\alpha\) is unsatisfiable.
\end{enumerate}
So \(\alpha\) is valid iff \(\alpha^{\mathrm d}\) is unsatisfiable, and
\(\alpha\) is unsatisfiable iff \(\alpha^{\mathrm d}\) is valid.

\emph{5. Finite connectives.}  For NNF presentations
\(\chi_1,\ldots,\chi_k\), the right-associated finite connectives of
\cref{sec:prop-syntax-semantics-01} satisfy
\[
 \Bigl(\bigwedge_{j=1}^{k}\chi_j\Bigr)^{\mathrm d}
   =\bigvee_{j=1}^{k}\chi_j^{\mathrm d},
 \qquad
 \Bigl(\bigvee_{j=1}^{k}\chi_j\Bigr)^{\mathrm d}
   =\bigwedge_{j=1}^{k}\chi_j^{\mathrm d}.
\]
For \(k=0\) these read \(\topv^{\mathrm d}=\bottom\) and
\(\bottom^{\mathrm d}=\topv\).  For \(k>0\), write \(C_i,D_i\) for the
conjunction and disjunction recursions of the \(\chi_j\), and \(C_i',D_i'\)
for those of the \(\chi_j^{\mathrm d}\).  Downward induction on \(i\) gives
\(C_i^{\mathrm d}=D_i'\): \(C_k^{\mathrm d}=\chi_k^{\mathrm d}=D_k'\), and for
\(i<k\),
\(C_i^{\mathrm d}=(\chi_i\land C_{i+1})^{\mathrm d}
 =(\chi_i^{\mathrm d}\lor C_{i+1}^{\mathrm d})
 =(\chi_i^{\mathrm d}\lor D_{i+1}')=D_i'\).  The same induction with the two
clauses exchanged gives \(D_i^{\mathrm d}=C_i'\).  Taking \(i=1\) proves the
display.  Since the dual of a literal is that literal, the dual of a finite
conjunction of literals is the finite disjunction of the same literals in the
same order, and conversely.  A second downward induction of the same shape,
using the truth clauses for \(\land\) and \(\lor\), shows that \(v\) satisfies
\(\bigwedge_{j}\chi_j\) iff it satisfies every \(\chi_j\), and satisfies
\(\bigvee_j\chi_j\) iff it satisfies some \(\chi_j\); for \(k=0\) this is the
value of \(\topv\) and of \(\bottom\).

\emph{6. Conjunctive normal forms go to disjunctive normal forms.}  A
conjunctive-normal-form formula is either \(\topv\), whose dual \(\bottom\) is
the empty disjunction and hence in disjunctive normal form, or the formula
\(N_1\) built from clauses \(C_1,\ldots,C_m\) with \(m>0\).  The recursion
defining \(N_m,\ldots,N_1\) is exactly the right-associated conjunction, so
\(N_1=\bigwedge_{i=1}^m\operatorname{form}_{\mathsf{CPL}}(C_i)\) and, by
part~5,
\(N_1^{\mathrm d}=\bigvee_{i=1}^m\operatorname{form}_{\mathsf{CPL}}(C_i)^{\mathrm d}\).
If \(C_i=\varnothing\), then
\(\operatorname{form}_{\mathsf{CPL}}(C_i)^{\mathrm d}=\bottom^{\mathrm d}=\topv\),
the empty conjunction.  Otherwise
\(\operatorname{form}_{\mathsf{CPL}}(C_i)\) is the right-associated
disjunction of the members of \(C_i\) in strictly increasing code order, and
its dual is their right-associated conjunction in the same order.  So
\(N_1^{\mathrm d}\) is a finite disjunction of finite conjunctions of literals:
it is in disjunctive normal form.

\emph{7. Disjunctive normal forms go to conjunctive normal forms up to
equivalence.}  Let \(\delta=\bigvee_{i=1}^n\kappa_i\) with
\(\kappa_i=\bigwedge_{j=1}^{k_i}L_{ij}\) be in disjunctive normal form.  If
\(n=0\), then \(\delta=\bottom\) and \(\delta^{\mathrm d}=\topv\), which is in
conjunctive normal form.  If \(n>0\), then by part~5
\(\delta^{\mathrm d}=\bigwedge_{i=1}^n\kappa_i^{\mathrm d}\) with
\(\kappa_i^{\mathrm d}=\bigvee_{j=1}^{k_i}L_{ij}\).  Put
\(C_i=\{L_{i1},\ldots,L_{ik_i}\}\), a clause, and let \(N_1\) be the
conjunctive-normal-form formula built from \(C_1,\ldots,C_n\).  For each \(i\)
and valuation \(v\): \(v\entails\kappa_i^{\mathrm d}\) iff \(v\) satisfies some
\(L_{ij}\) (part~5) iff \(v\) satisfies some member of \(C_i\) iff
\(v\entails\operatorname{form}_{\mathsf{CPL}}(C_i)\), the last step by
\cref{lem:prop-clause-satisfaction}(1); when \(k_i=0\) all of these fail.  By
part~5 again and \cref{lem:prop-clause-satisfaction}(2),
\(v\entails\delta^{\mathrm d}\) iff \(v\) satisfies every
\(\operatorname{form}_{\mathsf{CPL}}(C_i)\) iff \(v\entails N_1\).  Hence
\(\delta^{\mathrm d}\equiv N_1\).  When each \(\kappa_i\) lists distinct
literals in strictly increasing code order, \(\kappa_i^{\mathrm d}\) is the
string \(\operatorname{form}_{\mathsf{CPL}}(C_i)\) and \(\delta^{\mathrm d}\)
is the string \(N_1\).

\emph{8. The class inclusion in part~7 holds only up to equivalence.}  The
formula \((p_1\land p_1)\) is in disjunctive normal form: one disjunct, the
conjunction of the literal \(p_1\) listed twice.  Its dual \((p_1\lor p_1)\)
is not in conjunctive normal form.  It is not \(\topv\), which is a negation.
It is not \(N_1\) for \(m\geq 2\) clauses, since that is a conjunction and, by
Theorem~\ref{thm:prop-unique-readability}, no formula is both a conjunction
and a disjunction.  It is not \(\operatorname{form}_{\mathsf{CPL}}(C)\) for a
single clause \(C\): \(\operatorname{form}_{\mathsf{CPL}}(\varnothing)=\bottom\),
a one-member clause gives a literal, and for a clause of at least two members
unique readability would force the first two listed literals both to be
\(p_1\), whereas the members of a clause are distinct and listed in strictly
increasing code order.  It is logically equivalent to the conjunctive normal
form \(\operatorname{form}_{\mathsf{CPL}}(\{p_1\})=p_1\), as part~7 requires.

Parts~6--8 give the exact content of the final clause of the statement:
dualization maps every conjunctive normal form to a disjunctive normal form,
and maps every disjunctive normal form to a formula logically equivalent to the
conjunctive normal form with the same literal sets, which is that conjunctive
normal form itself whenever each disjunct lists distinct literals in
strictly increasing code order.
\end{proof}
\end{annexitem}

\begin{annexitem}{annex:p2.3.9}{thm:p02-interpolant-truth-table-correctness}{Correctness of the truth-table interpolant construction}
\begin{proof}
Let \(\theta\) be the output and let \(X\) be the common-variable set.  If a
valuation \(v\) satisfies \(\varphi\), its restriction \(v|_X\) is among the
selected assignments, so \(v\) satisfies the corresponding complete
conjunction and hence \(\theta\).  Thus
\(\varphi\entails\theta\).

Suppose instead that some valuation \(v\) satisfies \(\theta\) and falsifies
\(\psi\).  It satisfies a selected conjunction, whose assignment
\(u:X\to\{0,1\}\) therefore extends to a valuation \(w\) satisfying
\(\varphi\).  Amalgamate \(w\) on the variables occurring only in
\(\varphi\), \(v\) on those occurring only in \(\psi\), and their common
value \(u\) on \(X\).  The resulting valuation satisfies \(\varphi\) and
falsifies \(\psi\), contrary to the hypothesis.  Hence
\(\theta\entails\psi\).  By construction, every variable of \(\theta\) lies
in \(X\).  The empty selected family gives \(\bottom\), and an empty
conjunction gives \(\topv\), so the boundary cases obey the same argument.
\end{proof}
\end{annexitem}

\begin{annexitem}{annex:p2.3.12}{thm:p02-dual-horn-sat}{Dual-Horn-SAT tractability}
\begin{proof}
Complement every variable throughout the instance.  A clause with at most one
negative literal becomes a clause with at most one positive literal, and the
map is a linear-time bijection between satisfying assignments (complement the
assignment in the same way).  Apply the linear-time Horn algorithm
(Method~\ref{def:p02-w02-p02-019-horn-method}, with its \(O(N)\) bound
proved in Theorem~\ref{thm:p02-w02-p02-019-horn-procedure}) and map its
answer back.  Complementation rewrites each literal occurrence once, so the
whole procedure takes \(O(N)\) steps.
\end{proof}
\end{annexitem}

%% Assembling items for the decision and complexity statements of
%% sec:p02-prop-decision-complexity; the proofs are in
%% sec:annex-p02-w02-p02-018-decision-certificates.
\begin{annexitem}{annex:p02-w02-p02-018-decidability}{thm:p02-decidability}{Decidability of classical propositional logic}
\begin{proof}
Satisfiability, validity, logical equivalence and finite-premise consequence
are membership of the encoded input in \(\mathsf{SAT}\), \(\mathsf{TAUT}\),
\(\mathsf{EQV}\) and \(\mathsf{CONS}\) of
Convention~\ref{def:p02-w02-p02-018-input}.  By
Theorem~\ref{thm:p02-w02-p02-018-decision},
Method~\ref{def:p02-w02-p02-018-truth-table} halts on every bit string and
accepts exactly the members of each of these sets, within
\(c\,(2^k+1)\,m^2\) steps.  The scope sentence is
Theorem~\ref{thm:p02-w02-p02-018-limits}.
\end{proof}
\end{annexitem}

\begin{annexitem}{annex:p02-w02-p02-018-sat-np-complete}{thm:p02-sat-np-complete}{SAT NP-completeness}
\begin{proof}
Membership \(\mathsf{SAT}\in\mathsf{NP}\) is
Theorem~\ref{thm:p02-w02-p02-018-membership}.  Every \(L\in\mathsf{NP}\)
reduces to \(\mathsf{SAT}\) by
Theorem~\ref{thm:p02-w02-p02-018-cook-levin}.
\end{proof}
\end{annexitem}

\begin{annexitem}{annex:p02-w02-p02-018-taut-conp-complete}{thm:p02-taut-conp-complete}{TAUT is coNP-complete}
\begin{proof}
This is Corollary~\ref{thm:p02-w02-p02-018-taut-conp}.
\end{proof}
\end{annexitem}
